\documentclass[a4paper,10pt]{article}

\usepackage{amsmath,amssymb,amsthm}
\usepackage{braket}
\usepackage{graphicx}
\usepackage{indentfirst}
\usepackage[a4paper,margin=3cm]{geometry}
\usepackage{hyperref}

\hypersetup{
    colorlinks=true,
    linkcolor=blue,
    citecolor=blue,
    urlcolor=blue
}

\newcommand{\I}{\mathbb{I}}
\newcommand{\Tr}{\operatorname{Tr}}

\title{Polarimetry and Measurement Bases for Entangled Photons}
\author{M. Gil de Oliveira}
\date{\today}

\begin{document}

\maketitle

These notes develop polarization analyzers for measurements on entangled
photon pairs. The central task is to characterize the shared state at the
measurement stations and choose bases adapted to its correlations.
Wave-plate transformations provide the implementation; an ideal maximally
entangled state supplies a useful reference, and tomography determines
whether that description is appropriate for the apparatus.

\section{Wave plates}

In a birefringent material, orthogonal polarization components acquire
different phases as they propagate. Let $\ket{0}$ and $\ket{1}$ denote the
principal polarization axes of a wave plate. In this basis, its action is

\begin{equation}
    \overline{U}_{\phi_0,\phi_1}
    =
    \begin{pmatrix}
        e^{i\phi_0} & 0 \\
        0 & e^{i\phi_1}
    \end{pmatrix}.
    \label{eq:wave-plate-principal-basis}
\end{equation}

Suppose that the principal axes are rotated by an angle $\theta$ relative to
the horizontal and vertical directions, so that
$\ket{0}=R_\theta\ket{h}$ and $\ket{1}=R_\theta\ket{v}$. In the
$\{\ket{h},\ket{v}\}$ basis, the rotation operator is

\begin{equation}
    R_\theta
    =
    \begin{pmatrix}
        \cos\theta & -\sin\theta \\
        \sin\theta & \cos\theta
    \end{pmatrix},
    \qquad
    R_\theta^{-1}=R_\theta^\dagger=R_{-\theta}.
    \label{eq:rotation}
\end{equation}

The wave-plate operator in the horizontal--vertical basis is therefore

\begin{equation}
    U_{\phi_0,\phi_1}(\theta)
    =R_\theta\overline{U}_{\phi_0,\phi_1}R_{-\theta}
    =\frac{e^{i\phi_0}+e^{i\phi_1}}{2}\I
    +\frac{e^{i\phi_0}-e^{i\phi_1}}{2}S_{2\theta},
    \label{eq:general-wave-plate}
\end{equation}

where

\begin{equation}
    S_\theta
    =
    \begin{pmatrix}
        \cos\theta & \sin\theta \\
        \sin\theta & -\cos\theta
    \end{pmatrix}
    =R_\theta Z
    =ZR_{-\theta},
    \qquad
    Z=
    \begin{pmatrix}
        1 & 0 \\
        0 & -1
    \end{pmatrix}.
    \label{eq:S-operator}
\end{equation}

This decomposition gives the useful composition rule

\begin{equation}
    S_{\theta_2}S_{\theta_1}=R_{\theta_2-\theta_1}.
    \label{eq:S-composition}
\end{equation}

\subsection{Half-wave and quarter-wave plates}

For a half-wave plate (HWP), we choose $\phi_0=0$ and $\phi_1=\pi$.
Equation~\eqref{eq:general-wave-plate} then yields

\begin{equation}
    U^{\lambda/2}_\theta=S_{2\theta}.
    \label{eq:hwp}
\end{equation}

For a quarter-wave plate (QWP), we choose $\phi_0=-\pi/4$ and
$\phi_1=\pi/4$. With this global-phase convention,

\begin{equation}
    U^{\lambda/4}_\theta
    =\frac{1}{\sqrt{2}}\left(\I-iS_{2\theta}\right),
    \qquad
    \left(U^{\lambda/4}_\theta\right)^\dagger
    =\frac{1}{\sqrt{2}}\left(\I+iS_{2\theta}\right).
    \label{eq:qwp}
\end{equation}

\section{Quarter-wave-plate--half-wave-plate analyzer}

Consider an analyzer consisting of a QWP at angle $\theta_1$, followed by an
HWP at angle $\theta_2$ and a polarizing beam splitter (PBS). The wave plates
apply the unitary

\begin{equation}
    U_w=U^{\lambda/2}_{\theta_2}U^{\lambda/4}_{\theta_1}.
    \label{eq:analyzer-unitary}
\end{equation}

Detection at the horizontal output of the PBS realizes a projection onto the
state

\begin{align}
    \ket{\psi_{\theta_1,\theta_2}}
    &=U_w^\dagger\ket{h} \notag\\
    &=\frac{1}{\sqrt{2}}
      \left(\I+iS_{2\theta_1}\right)S_{2\theta_2}\ket{h} \notag\\
    &=\frac{1}{\sqrt{2}}
      \left[S_{2\theta_2}+iR_{2(\theta_1-\theta_2)}\right]\ket{h}.
    \label{eq:analyzer-state-compact}
\end{align}

Because $Z\ket{h}=\ket{h}$, Eq.~\eqref{eq:S-operator} implies
$S_\theta\ket{h}=R_\theta\ket{h}$. Hence the analyzer state can be written
explicitly as

\begin{equation}
    \boxed{
    \begin{aligned}
        \ket{\psi_{\theta_1,\theta_2}}
        =\frac{1}{\sqrt{2}}\big[&
        \left(\cos 2\theta_2
        +i\cos 2(\theta_1-\theta_2)\right)\ket{h}\\
        {}+&\left(\sin 2\theta_2
        +i\sin 2(\theta_1-\theta_2)\right)\ket{v}\big].
    \end{aligned}}
    \label{eq:analyzer-state}
\end{equation}

\subsection{Standard polarimetry settings}

A polarization measurement is obtained by recording optical power or photon
counts at the PBS outputs. Varying $\theta_1$ and $\theta_2$ selects the six states commonly
used for polarization tomography. Define

\begin{equation}
    \ket{d}=\frac{\ket{h}+\ket{v}}{\sqrt{2}},
    \quad
    \ket{a}=\frac{\ket{h}-\ket{v}}{\sqrt{2}},
    \quad
    \ket{r}=\frac{\ket{h}-i\ket{v}}{\sqrt{2}},
    \quad
    \ket{l}=\frac{\ket{h}+i\ket{v}}{\sqrt{2}}.
    \label{eq:polarization-states}
\end{equation}

With the phase convention of Eq.~\eqref{eq:qwp}, convenient analyzer settings
are

\begin{alignat}{3}
    \ket{\psi_{0,0}}             &=e^{i\pi/4}\ket{h},
    &\qquad
    \ket{\psi_{0,\pi/4}}        &=e^{-i\pi/4}\ket{v}, \notag\\
    \ket{\psi_{\pi/4,\pi/8}}   &=e^{i\pi/4}\ket{d},
    &
    \ket{\psi_{\pi/4,3\pi/8}} &=-e^{-i\pi/4}\ket{a}, \notag\\
    \ket{\psi_{0,\pi/8}}       &=e^{i\pi/4}\ket{r},
    &
    \ket{\psi_{0,3\pi/8}}      &=e^{i\pi/4}\ket{l}.
    \label{eq:tomography-settings}
\end{alignat}

The phase factors in Eq.~\eqref{eq:tomography-settings} do not affect the
corresponding projectors or measured powers.

\section{The shared polarization state}

The starting point of a Bell experiment is the two-photon state delivered to
the measurement stations. A fiber between the source and one station may
change the polarization bases without reducing the entanglement. It is
therefore useful to characterize the delivered state first and then select
measurement bases adapted to its correlations. Identifying which part of a
basis mismatch originated in the source and which in propagation is not
necessary for this purpose.

We label the local and traveling photons $L$ and $T$, respectively, and use
the same horizontal--vertical convention in each arm. The reference Bell
state is
\begin{equation}
    \ket{\Phi^+}
    =\frac{\ket h_L\ket h_T+\ket v_L\ket v_T}{\sqrt{2}}.
    \label{eq:spdc-bell-state}
\end{equation}
This is a reference for describing the correlations, rather than an
assumption about the state produced by the apparatus.

\subsection{Maximally entangled states and relative polarization bases}

Every pure maximally entangled two-qubit state can be written as
$(A\otimes B)\ket{\Phi^+}$ for suitable local unitaries. The identity
\begin{equation}
    (A\otimes\I)\ket{\Phi^+}
    =(\I\otimes A^{\mathsf T})\ket{\Phi^+}
    \label{eq:bell-transpose-identity}
\end{equation}
then gives the equivalent representation
\begin{equation}
    \boxed{\ket{\Psi_U}=(\I\otimes U)\ket{\Phi^+}},
    \qquad U=BA^{\mathsf T}.
    \label{eq:shared-state}
\end{equation}
All transposes and complex conjugates in these notes refer to the $h/v$
basis. The unitary $U$ is defined up to an irrelevant overall phase and
describes the relative polarization bases of the pair. If the source has
relative transformation $U_s$ and the traveling photon crosses a unitary
fiber $U_f$, the delivered state has $U=U_fU_s$. An additional local
transformation $U_L$ changes this to $U=U_fU_sU_L^{\mathsf T}$.

This description is exact for pure maximally entangled states. A measured
mixed or nonmaximally entangled state should instead be represented by its
density matrix $\rho_{LT}$; it must not be forced into
Eq.~\eqref{eq:shared-state} without checking the approximation.

\subsection{Perfectly correlated bases}

The reference state obeys
\begin{equation}
    (X\otimes X)\ket{\Phi^+}=\ket{\Phi^+},
    \qquad
    (Z\otimes Z)\ket{\Phi^+}=\ket{\Phi^+}.
\end{equation}
Consequently, the state in Eq.~\eqref{eq:shared-state} satisfies
\begin{equation}
    \boxed{
    \begin{aligned}
        (X\otimes UXU^\dagger)\ket{\Psi_U}&=\ket{\Psi_U},\\
        (Z\otimes UZU^\dagger)\ket{\Psi_U}&=\ket{\Psi_U}.
    \end{aligned}}
    \label{eq:matched-axes}
\end{equation}
Thus a local $H/V$ measurement is perfectly correlated with the remote
basis $\{U\ket h,U\ket v\}$, and a local $D/A$ measurement is perfectly
correlated with $\{U\ket d,U\ket a\}$. The remote states need not be linear
in the laboratory frame.

These are examples of a continuous family of correlated bases. For any
qubit polarization observable $A=\boldsymbol{n}\cdot\boldsymbol{\sigma}$
with unit vector $\boldsymbol n$,
\begin{equation}
    (A\otimes UA^{\mathsf T}U^\dagger)\ket{\Psi_U}
    =\ket{\Psi_U}.
    \label{eq:general-matched-observable}
\end{equation}
Equivalently, selecting $\ket b$ locally prepares $U\ket{b^*}$ remotely.
The conjugation is immaterial for linear states, whose amplitudes can be
chosen real, but is essential for elliptical states. In particular,
$\ket{\Phi^+}$ is anticorrelated in $Y/Y$, since $Y^{\mathsf T}=-Y$.
There is no unique preferred pair of measurement bases: choosing axes at
one station fixes the corresponding axes at the other.

\section{Correlations and Bell-test measurement settings}

In these notes, a Bell test means separate local polarization measurements
whose correlations enter a Bell inequality. A Bell-state measurement,
which jointly distinguishes states in an entangled basis, is a different
operation.

For observables $A_i$ and $B_j$ with outcomes $\pm1$, define the
Clauser--Horne--Shimony--Holt (CHSH) combination
\begin{equation}
    S=E_{00}+E_{01}+E_{10}-E_{11},
    \qquad E_{ij}=\Tr\!\left[\rho_{LT}(A_i\otimes B_j)\right].
    \label{eq:chsh}
\end{equation}
The local hidden-variable bound is $|S|\leq2$. Perfect correlations along
two matching axes are useful for calibration but are not themselves the
four settings of an optimal CHSH test.

For $\ket{\Psi_U}$, one convenient choice is
\begin{equation}
    \begin{aligned}
        A_0&=Z,& A_1&=X,\\
        B_0&=U\frac{Z+X}{\sqrt{2}}U^\dagger,
        &B_1&=U\frac{Z-X}{\sqrt{2}}U^\dagger.
    \end{aligned}
    \label{eq:ideal-chsh-settings}
\end{equation}
These settings give $S=2\sqrt{2}$. Before the transformation $U$, the
positive-eigenvalue remote states are linear polarizations at $22.5^\circ$
and $-22.5^\circ$ relative to $h$. The corresponding directions on the
Poincar\'e sphere are the bisectors of the $X$ and $Z$ axes.

If it is preferable to retain linear analyzers remotely, the basis
adaptation can instead be placed locally. For example,
\begin{equation}
    \begin{aligned}
        A_0&=U^{\mathsf T}ZU^*,&
        A_1&=U^{\mathsf T}XU^*,\\
        B_0&=\frac{Z+X}{\sqrt{2}},&
        B_1&=\frac{Z-X}{\sqrt{2}}
    \end{aligned}
    \label{eq:local-adapted-chsh}
\end{equation}
gives the same ideal CHSH value. Each station's settings depend on the
shared state and its own measurement choice, not on the other station's
choice in a particular trial.

\section{Characterizing the delivered state}

\subsection{Tomography after the delivery fiber}

For an SPDC source, the proposed reference plane is the output of the
delivery fiber, with the input coupling and fiber routing left in their
operational configuration. Light emerging from the fiber is collimated
into a temporary polarization analyzer, and two-photon tomography is
performed using coincidences with the other arm. This characterizes the
state in the spatial mode actually delivered to that plane.

Putting the diagnostic analyzer after the fiber avoids modifying the
input collection when inserting or removing polarization optics.
Changing the input coupling can change both the collected spatial mode
and, if the fiber is disturbed, its polarization transformation.
Tomography at this plane does not include optics farther downstream,
such as a telescope; those must be included in a later calibration if
they change the polarization.

Full two-qubit polarization tomography reconstructs $\rho_{LT}$ without
assuming a Bell state~\cite{James2001Measurement}. A convenient complete
measurement set is the nine combinations of $H/V$, $D/A$, and $R/L$ at
the two stations, recording both PBS outputs. With
$\sigma_0=\I$ and $(\sigma_x,\sigma_y,\sigma_z)=(X,Y,Z)$,
\begin{equation}
    \rho_{LT}
    =\frac14\sum_{\mu,\nu=0,x,y,z}
        t_{\mu\nu}\,\sigma_\mu\otimes\sigma_\nu,
    \qquad
    t_{\mu\nu}=\Tr\!\left[
        \rho_{LT}(\sigma_\mu\otimes\sigma_\nu)
    \right],
    \label{eq:two-photon-tomography}
\end{equation}
where $t_{00}=1$ leaves fifteen independent real parameters. The measured
marginals supply the single-photon terms as well as the joint correlations.

The reconstruction should account for detector efficiencies, unequal
acquisition times, and accidental coincidences. Source brightness and
fiber polarization must remain sufficiently stable during the acquisition.
A physical density-matrix fit and statistical resampling of the counts
can be used to estimate the uncertainties of derived quantities.
These are measurements of the detected polarization ensemble; collection
biases must be included when interpreting the result.

\subsection{Entanglement and alignment with a reference Bell state}

Concurrence assesses entanglement independently of local basis
rotations~\cite{Wootters1998Concurrence}. For a normalized pure state,
\begin{equation}
    \ket\psi=a\ket{hh}+b\ket{hv}+c\ket{vh}+d\ket{vv},
    \qquad
    M=\begin{pmatrix}a&b\\c&d\end{pmatrix},
\end{equation}
it is
\begin{equation}
    C(\ket\psi)=2|\det M|=2|ad-bc|.
\end{equation}
For a general reconstructed density matrix, define
$\widetilde\rho=(Y\otimes Y)\rho_{LT}^*(Y\otimes Y)$. If
$\lambda_1\geq\lambda_2\geq\lambda_3\geq\lambda_4$ are the square roots
of the eigenvalues of $\rho_{LT}\widetilde\rho$, then
\begin{equation}
    C(\rho_{LT})
    =\max(0,\lambda_1-\lambda_2-\lambda_3-\lambda_4).
    \label{eq:concurrence}
\end{equation}
Concurrence is not the determinant of the density matrix. It is zero for
separable two-qubit states and one for maximally entangled states.
An unknown stable local unitary does not change it.

A concurrence close to one supports a nearly maximally entangled
description but does not determine the measurement bases. To quantify
how well a local controller could align the state with the reference,
one can also evaluate
\begin{equation}
    F_{\mathrm{opt}}
    =\max_V
    \bra{\Phi^+}(V\otimes\I)\rho_{LT}
        (V^\dagger\otimes\I)\ket{\Phi^+},
    \label{eq:optimized-bell-fidelity}
\end{equation}
where the maximization is over local unitaries. The maximizing $V$ gives
a candidate correction; for Eq.~\eqref{eq:shared-state}, it is $U^*$ and
$F_{\mathrm{opt}}=1$. A unitary cannot remove the residual mixedness of
an imperfect state. Both concurrence and optimized fidelity should be
reported with their statistical uncertainties.

\subsection{Choosing settings from an imperfect state}

When the measured state is not adequately described by a pure maximally
entangled state, its correlations can be used directly. Let $T$ be the
$3\times3$ submatrix with entries $T_{ij}=t_{ij}$ for $i,j=x,y,z$.
For unit measurement directions $\boldsymbol a$ and $\boldsymbol b$,
\begin{equation}
    E(\boldsymbol a,\boldsymbol b)
    =\boldsymbol a^{\mathsf T}T\boldsymbol b.
    \label{eq:measured-correlations}
\end{equation}
Substituting this expression into Eq.~\eqref{eq:chsh} allows the four
directions to be optimized using the reconstructed state. Hardware
restrictions must be imposed in that optimization: a station with only
an HWP and PBS has linear directions
$\boldsymbol b=(\sin2\beta,0,\cos2\beta)$, whereas a QWP--HWP analyzer can
select any direction. This procedure neither presumes a Bell state nor
requires separating source imperfections from propagation effects.

\section{Implementing the selected measurement bases}

\subsection{From a target projector to wave-plate angles}

Once the target state $\ket b=b_h\ket h+b_v\ket v$ at an analyzer input
has been chosen, the required wave-plate transformation satisfies
\begin{equation}
    \boxed{
        U_w^\dagger\Pi_hU_w=\ket b\bra b,
        \qquad \Pi_h=\ket h\bra h
    }.
    \label{eq:target-projector}
\end{equation}
The vertical PBS port automatically realizes
$\I-\ket b\bra b$, so a single pair of angles specifies the whole
orthonormal measurement basis.

For a normalized state with amplitudes $b_h,b_v$, define
\begin{equation}
    c_x=2\operatorname{Re}(b_h^*b_v),\qquad
    c_y=2\operatorname{Im}(b_h^*b_v),\qquad
    c_z=|b_h|^2-|b_v|^2.
    \label{eq:stokes-definition}
\end{equation}
The projector is $(\I+c_xX+c_yY+c_zZ)/2$. From
Eq.~\eqref{eq:analyzer-state}, the analyzer selects the Stokes vector
\begin{equation}
    \begin{aligned}
        s_x&=\sin(2\theta_1)\cos(4\theta_2-2\theta_1),\\
        s_y&=-\sin(4\theta_2-2\theta_1),\\
        s_z&=\cos(2\theta_1)\cos(4\theta_2-2\theta_1).
    \end{aligned}
    \label{eq:analyzer-stokes}
\end{equation}
Setting $\boldsymbol s=\boldsymbol c$ and choosing
$\cos(4\theta_2-2\theta_1)=\sqrt{c_x^2+c_z^2}\geq0$ gives
\begin{equation}
    \boxed{
    \begin{aligned}
        \theta_1&=\frac12\operatorname{atan2}(c_x,c_z),\\[3pt]
        \theta_2&=\frac14\left[
            \operatorname{atan2}(c_x,c_z)
            +\operatorname{atan2}\!\left(-c_y,\sqrt{c_x^2+c_z^2}\right)
        \right].
    \end{aligned}}
    \label{eq:compensation-angles}
\end{equation}
Here $\operatorname{atan2}(y,x)$ is the polar angle
$\gamma\in(-\pi,\pi]$ of $(x,y)$, defined for $(x,y)\neq(0,0)$ by
\begin{equation}
    \cos\gamma=\frac{x}{\sqrt{x^2+y^2}},
    \qquad
    \sin\gamma=\frac{y}{\sqrt{x^2+y^2}}.
\end{equation}
Unlike $\arctan(y/x)$, it keeps the signs of both coordinates, chooses the
correct quadrant, and handles $x=0$ when $y\neq0$.
At $c_x=c_z=0$, one may choose $\theta_1$ freely and use
\begin{equation}
    \theta_2=\frac{\gamma+2\theta_1}{4},
    \qquad \gamma=\operatorname{atan2}(-c_y,0).
\end{equation}

For the target linear state
$\ket{\alpha}=\cos\alpha\ket h+\sin\alpha\ket v$, one convenient solution
is $\theta_1=\alpha$, $\theta_2=\alpha/2$, with equivalent branches
understood. In particular, $H/V$ uses $(0,0)$ and $D/A$ uses
$(45^\circ,22.5^\circ)$.

\subsection{Choosing angle branches to reduce analyzer movement}

The inverse in Eq.~\eqref{eq:compensation-angles} is not unique. Its
choice of a nonnegative cosine is convenient algebraically, but need not
minimize movement between successive projectors. An alternative convention
places both optical angles in $[0,\pi/2)$ while retaining the sign of the
projection onto the selected QWP direction. For a unit Stokes vector,
define
\begin{equation}
    \boxed{
    \begin{aligned}
        q&=\left[\frac12\operatorname{atan2}(c_x,c_z)\right]
            \bmod\frac{\pi}{2},\\
        t&=c_x\sin(2q)+c_z\cos(2q),\\
        h&=\left[\frac{2q+\operatorname{atan2}(-c_y,t)}{4}\right]
            \bmod\frac{\pi}{2},
        \qquad (\theta_1,\theta_2)=(q,h).
    \end{aligned}}
    \label{eq:movement-angle-convention}
\end{equation}
Here modulo returns the representative in $[0,\pi/2)$. At the circular
poles $c_x=c_z=0$, set $q=0$. Away from the poles, the chosen QWP direction
is parallel or antiparallel to $(c_x,c_z)$, so
$t=\pm\sqrt{c_x^2+c_z^2}$. The signed value $t$, rather than its absolute
value, selects the required HWP branch. Substitution into
Eq.~\eqref{eq:analyzer-stokes} gives the requested Stokes vector.
The reduction of $q$ modulo $\pi/2$ is accompanied by this recalculation
of $h$; it is not a claim that a QWP alone has that angular period.

With the circular-state definitions of Eq.~\eqref{eq:polarization-states},
$R$ has $c_y=-1$ and $L$ has $c_y=+1$. Equation~\eqref{eq:movement-angle-convention}
gives the following sequence, with all angles expressed in degrees:
\begin{center}
\begin{tabular}{c c c l}
Projector & QWP $q$ & HWP $h$ & Movement from previous setting\\
\hline
$H$ & $0$ & $0$ & initial setting\\
$R$ & $0$ & $22.5$ & HWP $+22.5^\circ$\\
$V$ & $0$ & $45$ & HWP $+22.5^\circ$\\
$L$ & $0$ & $67.5$ & HWP $+22.5^\circ$\\
$A$ & $45$ & $67.5$ & QWP $+45^\circ$\\
$D$ & $45$ & $22.5$ & HWP $-45^\circ$
\end{tabular}
\end{center}
Every transition moves exactly one plate. The five movements sum to
$157.5^\circ$: $45^\circ$ on the QWP and $112.5^\circ$ on the HWP.

The ordering can be obtained by a small discrete optimization. Assign
one vertex to each of the six angle pairs in the table and connect two
vertices only when one plate angle is unchanged. Give each edge the
absolute change of the other angle. Enumerate the $6!=720$ projector
orders, discard those containing a forbidden transition, and minimize
the sum of the five edge costs. This yields the displayed order (or its
reverse), with cost $157.5^\circ$. Five movements are also the minimum
possible number for six distinct projectors. The angular minimum stated
here is for these fixed representatives and an open sequence: it excludes
initial positioning and returning to the start, and is not an unrestricted
optimization over every equivalent angle branch. For arbitrary target
projectors, Eq.~\eqref{eq:movement-angle-convention} is a deterministic
convention, not a trajectory optimizer.

If one PBS output per station is sampled sequentially, the resulting
$36$ joint projector measurements can use a serpentine traversal. Let
$P=(H,R,V,L,A,D)$. Advance the first station through $P$; at successive
settings, traverse $P$ at the second station alternately forwards and
backwards. At a row boundary, the second station remains at its last
setting while only one plate at the first station moves. Thus all
$36$ pairs require only $35$ individual plate movements between
measurements, the minimum possible count. Their total angular travel is
$6(157.5^\circ)+157.5^\circ=1102.5^\circ$. This construction establishes
a minimum movement count, not a globally minimum angular path through
all joint settings. If both PBS outputs are recorded, the nine joint
basis settings discussed earlier already supply the $36$ outcome
combinations; the sequential-projector schedule is then unnecessary.

These angles refer to optical axes; fixed calibrated mount offsets must
be added when commanding the motors. Unchanged plate targets should not
be commanded again. The travel calculation assumes the indicated signed
rotations are available and excludes homing and settling delays. Minimum
angular travel need not imply minimum acquisition time when motor speeds,
direction-dependent behavior, or settling requirements differ. Changes
of branch near coordinate boundaries also mean that this convention does
not guarantee smooth motion for a continuously varying target state.

\subsection{Adapting either station or applying a fixed correction}

For $\ket{\Psi_U}$, ideal reference projectors $\Pi_a$ and $\Pi_b$ can
be implemented by leaving the local projector unchanged and adapting
the remote one to $U\Pi_bU^\dagger$. Alternatively, the remote
projector can remain unchanged while the local one becomes
\begin{equation}
    \widetilde{\Pi}_a=U^{\mathsf T}\Pi_aU^*.
    \label{eq:local-effective-projector}
\end{equation}
Both choices reproduce the corresponding joint probabilities of
$\ket{\Phi^+}$. The analyzer target kets are respectively $U\ket b$
or $U^{\mathsf T}\ket a$; their Stokes vectors are inserted into
Eq.~\eqref{eq:compensation-angles}. Two adjustable plates suffice for
each projector even though $U$ itself has three parameters.

A different implementation restores the reference state before analysis:
\begin{equation}
    \boxed{(U^*\otimes U)\ket{\Phi^+}=\ket{\Phi^+}.}
    \label{eq:local-channel-correction}
\end{equation}
A general local controller, such as a QWP--HWP--QWP sequence, can apply
$V=U^*$. Independent analyzers may then use the usual reference settings.
If the same controller also performs local analysis, its setting must
change with the local basis: an analyzer $A_a$ satisfying
$A_a^\dagger\ket h=\ket a$ requires total transformation $A_aU^*$.
The existence of a fixed correction therefore does not remove the need
to select local measurement bases.

For a real rotation $U=R_\beta$, the local correction is the same
rotation $R_\beta$. For $U=\operatorname{diag}(1,e^{i\delta})$, it is
$\operatorname{diag}(1,e^{-i\delta})$. This restores joint correlations;
it does not reverse a channel on an arbitrary independent photon.
The reduced state of each photon remains $\I/2$ in the ideal model,
so calibration information comes from coincidences rather than singles.

\subsection{Horizontal--vertical channel compensation as an example}

The earlier channel picture is recovered when a specified input state
passes through a known unitary $U_c$. For an input density matrix $\rho$,
the horizontal detection probability is
\begin{equation}
    p(h)=\Tr\!\left[
        U_c^\dagger U_w^\dagger\Pi_hU_wU_c\,\rho
    \right].
\end{equation}
To measure the input projector $\Pi_b$, choose
\begin{equation}
    U_w^\dagger\Pi_hU_w=U_c\Pi_bU_c^\dagger.
    \label{eq:arbitrary-basis-condition}
\end{equation}
For $H/V$, the target is simply $U_c\ket h$. The other PBS output
then measures the input $v$ projector by orthogonality.

This one-basis calibration leaves a relative phase undetermined:
\begin{equation}
    U_wU_c=e^{i\chi}
    \begin{pmatrix}1&0\\0&e^{i\delta}\end{pmatrix}.
    \label{eq:residual-phase}
\end{equation}
An independent $D/A$ analyzer after this transformation measures,
at the input, $(\ket h\pm e^{-i\delta}\ket v)/\sqrt2$. Hence a setting
that restores $H/V$ need not restore $D/A$.
Rotating the HWP inside the QWP--HWP pair by $22.5^\circ$ also does
not generally select the desired $D/A$ basis. Even without a channel,
the settings $(0,0)$ select $h$, while $(0,22.5^\circ)$ select $r$.
The appropriate target for an input diagonal measurement is $U_c\ket d$.
This missing relative phase is the same information supplied by a
second correlated axis in the shared-state description.

\subsection{Heralded calibration of individual bases}

When the maximally entangled model is justified, individual analyzer
targets can be calibrated directly through conditional tomography:
\begin{equation}
    \begin{aligned}
        {}_L\!\bra a\Psi_U\rangle
            &=\frac{1}{\sqrt2}U\ket{a^*}_T,\\
        {}_T\!\bra b\Psi_U\rangle
            &=\frac{1}{\sqrt2}U^{\mathsf T}\ket{b^*}_L.
    \end{aligned}
    \label{eq:heralded-targets}
\end{equation}
Each ideal herald outcome has probability $1/2$. For linear target states
the conjugation disappears. Heralding $h$ and $d$ locally therefore
provides the two remote targets for matched $H/V$ and $D/A$ measurements.
Conversely, heralding those states remotely provides the corresponding
local targets while retaining linear remote analyzers.

Tomography of each conditional state still needs all three Stokes
components, because a linear herald can prepare an elliptical state.
Once a target is reconstructed, Eq.~\eqref{eq:compensation-angles}
provides its analyzer settings. Recording both herald PBS outputs
also permits calibration and checks of the complementary target.

For a general measured state, the conditional remote density matrix is
instead
\begin{equation}
    \rho_{T|a}
    =\frac{
        \Tr_L[(\Pi_a\otimes\I)\rho_{LT}(\Pi_a\otimes\I)]
    }{
        \Tr[(\Pi_a\otimes\I)\rho_{LT}]
    }.
\end{equation}
It may be mixed, and complementary heralds need not produce orthogonal
states. The pure-state prescription cannot then be interpreted as an
exact restoration of all correlations; settings should be selected
from the measured state as described above.

\section{Measurement resources and reduced calibration}

\subsection{Complete tomography versus verification}

A general controller or QWP--HWP pair at one station and only an HWP at
the other do not provide complete two-photon tomography. The latter
station accesses only the span of $\I,X,Z$. In
Eq.~\eqref{eq:two-photon-tomography}, the four coefficients
$t_{0y},t_{xy},t_{yy},t_{zy}$ are therefore unobserved. Further remote
linear settings do not supply them. A temporary manual QWP at that
station supplies circular analysis; full initial tomography does not
require every wave plate to be motorized.

Incomplete measurements can nevertheless verify a chosen reference
state. For any physical $\rho_{LT}$,
\begin{equation}
    \bra{\Phi^+}\rho_{LT}\ket{\Phi^+}
    \geq\frac{\langle X\otimes X\rangle+
                  \langle Z\otimes Z\rangle}{2}.
    \label{eq:bell-fidelity-bound}
\end{equation}
This follows from the operator inequality
$\ket{\Phi^+}\bra{\Phi^+}\geq(X\otimes X+Z\otimes Z)/2$.
It can certify high fidelity after basis alignment without presuming a
pure state, but does not reconstruct a general density matrix or
separate source errors from channel errors.

\subsection{Reduced characterization under the ideal model}

Characterizing a process through correlations with another photon is
the principle of ancilla-assisted process tomography
\cite{Altepeter2003Ancilla}. Once the pure maximally entangled description
has been established to adequate accuracy, the following reduced
procedure determines its relative unitary $U$.

Measure $H/V$ remotely and perform full local tomography conditioned on
the $h$ outcome; repeat with remote $D/A$, conditioned on $d$. By
Eq.~\eqref{eq:heralded-targets}, the local states are
\begin{equation}
    \rho_{L|h}=U^{\mathsf T}\Pi_hU^*,
    \qquad
    \rho_{L|d}=U^{\mathsf T}\Pi_dU^*.
    \label{eq:two-remote-conditions}
\end{equation}
These determine $U$ up to global phase. Indeed, if two unitaries $W_1,W_2$
representing $U^{\mathsf T}$ reproduce both projectors, then
$W_2^\dagger W_1$ preserves both $\Pi_h$ and $\Pi_d$. The first condition
makes it diagonal in $h/v$ and the second forces its two phases to agree.
Thus six combinations of three local and two remote bases suffice in
this restricted model.

This is a calibration of the relative bases, not independent evidence
that the source is a Bell state or that the channel is unitary.
Calibrated local projectors are required to identify a numerical $U$;
otherwise coincidence feedback can locate useful controller settings
without identifying its Jones matrix. Drift during calibration and
departures from the assumed state must be checked with additional data.

\subsection{Choosing which elements must move}

For sequential Bell-test settings, a local QWP--HWP pair can implement
the adapted local bases of Eq.~\eqref{eq:local-adapted-chsh}, while a
single remote HWP selects linear bases. Both local plates may need to
move between local settings. A three-parameter fiber controller can
serve the same purpose if it realizes arbitrary unitaries and is
calibrated. Keeping it at a fixed correction setting requires a separate
analyzer if the local basis is to vary independently.

A remote QWP may remain installed for occasional full tomography
without needing its own motor. For this option, the order
\begin{equation}
    \text{rotating HWP}\ \longrightarrow\
    \text{manual QWP}\ \longrightarrow\ \text{PBS}
\end{equation}
is useful. With the QWP axes aligned to $h/v$, it commutes with the PBS
projectors and does not affect detection probabilities. The HWP alone
then selects linear angle $\alpha=2\theta_H$. Rotating the manual QWP
allows circular analysis during tomography. This reversed order has
different general angle formulas from the QWP--HWP analyzer derived
earlier; those formulas must not be used interchangeably.

If two bases must be sampled simultaneously with a beam splitter and two
PBS arms, independent QWP--HWP analyzers provide a direct implementation
of the two selected projectors. Any unitary polarization transformations
specific to the splitter paths should be included in the arm calibrations.
A common full correction followed by separate linear analyzers is another
option. Merely aligning one basis with a common two-plate pair leaves
the relative-phase ambiguity of Eq.~\eqref{eq:residual-phase}.
Neither configuration reverses depolarization or general
polarization-dependent loss.

The operational sequence is therefore to characterize the delivered
state with the collection geometry fixed, assess its entanglement and
uncertainty, select accessible Bell-test settings, and calibrate the
analyzers for those projectors. Correlation checks can subsequently
monitor stability without repeating full tomography at every step.

\bibliographystyle{unsrt}
\bibliography{../../refs}

\end{document}
